% !TeX program = lualatex
% halo:begin
% {
%   "version": 1,
%   "publish": true,
%   "course": "mathematical-analysis-i",
%   "section": "2.4",
%   "title": "数学分析（一）2.4：数列收敛的判别",
%   "categories": ["study-notes"],
%   "tags": ["mathematics"],
%   "excerpt": "单调有界定理及其应用、子列与收敛判别。"
% }
% halo:end
% 来源：IMG_1581.HEIC—IMG_1584.HEIC；子列部分见 IMG_1589.HEIC。
\RequirePackage{iftex}
\RequireLuaTeX
\documentclass[UTF8, fontset=fandol, a4paper, 12pt]{ctexart}
\usepackage[margin=2.5cm]{geometry}
\usepackage{amsmath,amssymb,amsthm}
\usepackage{enumitem,fancyhdr}
\usepackage[hidelinks]{hyperref}
\newcommand{\R}{\mathbb R}
\newcommand{\Nplus}{\mathbb N_{+}}
\theoremstyle{definition}
\newtheorem{theorem}{定理}
\renewcommand{\thetheorem}{2.\arabic{theorem}}
\newtheorem*{example}{例题}
\renewcommand{\proofname}{证明}
\setlist[enumerate]{leftmargin=2em,itemsep=0.2em,topsep=0.3em}
\setlength{\headheight}{16pt}
\pagestyle{fancy}
\fancyhf{}
\fancyhead[L]{\small 数学分析（一）}
\fancyhead[R]{\small 2.4\quad 数列收敛的判别}
\fancyfoot[C]{\thepage}
\renewcommand{\headrulewidth}{0.3pt}
\raggedbottom
\title{数列收敛的判别}
\author{Yuchen Zhang}
\date{整理日期：2026年9月29日}
\makeatletter
\renewcommand{\maketitle}{\begin{center}
{\LARGE\bfseries \@title\par}\vspace{0.6em}
\@author\par\vspace{0.2em}{\small\@date\par}
\end{center}}
\makeatother
\begin{document}
\maketitle
\setcounter{section}{3}
\setcounter{theorem}{8}
\section{数列收敛的判别}

\subsection*{单调有界定理}
\begin{theorem}[单调有界定理]
单调递增且有上界的实数列必收敛，其极限是该数列的上确界；
单调递减且有下界的实数列也必收敛，其极限是该数列的下确界。
\end{theorem}
\begin{proof}
设 $\{a_n\}$ 单调递增且有上界，令 $L=\sup\{a_n:n\in\Nplus\}$。
对任意 $\varepsilon>0$，由上确界的性质，存在 $n_0$ 使
$a_{n_0}>L-\varepsilon$。当 $n\ge n_0$ 时，
\[
L-\varepsilon<a_{n_0}\le a_n\le L,
\]
故 $|a_n-L|<\varepsilon$，即 $a_n\to L$。
对单调递减且有下界的数列，向 $\{-a_n\}$ 应用上述结论即可。
\end{proof}

\begin{example}[$p$ 级数的敛散性]
设 $s_n=\sum_{k=1}^{n}k^{-\alpha}$，其中 $\alpha\in\R$。
当 $\alpha>1$ 时，$\{s_n\}$ 收敛；当 $\alpha\le1$ 时，$s_n\to+\infty$。
\end{example}
\begin{proof}
$s_n$ 严格递增。若 $\alpha>1$，把正整数按
$[2^j,2^{j+1}-1]$ 分组，则第 $j$ 组之和不超过
$2^j(2^j)^{-\alpha}=2^{-j(\alpha-1)}$。
令 $q=2^{1-\alpha}<1$，由有限等比和公式，任一部分和都不超过
$1/(1-q)$，故 $\{s_n\}$ 有上界，因而收敛。
若 $\alpha\le1$，则 $k^{-\alpha}\ge k^{-1}$。调和级数的第 $j$ 组
$\sum_{k=2^j}^{2^{j+1}-1}k^{-1}\ge\frac12$，故其部分和无上界，
从而 $s_n\to+\infty$。
特别地，当 $\alpha=2$ 时，还可直接估计
\[
1+\sum_{k=2}^{n}\frac1{k^2}
<1+\sum_{k=2}^{n}\frac1{k(k-1)}
=2-\frac1n<2\qquad(n\ge2).
\]
\end{proof}

\newpage
\begin{example}[递推根式]
令 $a_1=\sqrt2$，$a_{n+1}=\sqrt{2+a_n}$。求 $\lim_{n\to\infty}a_n$。
\end{example}
\begin{proof}
由归纳法，$0<a_n<2$：初始项满足此式，且 $0<a_n<2$ 时
$0<a_{n+1}=\sqrt{2+a_n}<2$。又
\[
a_{n+1}>a_n
\quad\Longleftrightarrow\quad
2+a_n>a_n^2
\quad\Longleftrightarrow\quad
(2-a_n)(a_n+1)>0.
\]
故数列递增且有上界，设其极限为 $A$。由平方根的极限性质，
$A=\sqrt{2+A}$，故 $A\in\{2,-1\}$。由于 $A\ge a_1=\sqrt2>0$，所以 $A=2$。
\end{proof}

\begin{example}[从集合中取趋于上确界的递增数列]
设 $S\subset\R$ 非空、有上界，$a=\sup S$，且 $a\notin S$。
则存在严格递增的 $\{x_n\}\subset S$，使 $x_n\to a$。
\end{example}
\begin{proof}
先取 $x_1\in S$，使 $a-1<x_1<a$。已取 $x_n$ 后，因
$\max\{x_n,a-1/(n+1)\}<a$，由上确界的性质可取 $x_{n+1}\in S$，使
\[
x_{n+1}>\max\left\{x_n,a-\frac1{n+1}\right\}.
\]
于是 $\{x_n\}$ 严格递增，且 $a-1/n<x_n<a$ 对一切 $n$ 成立，
故 $x_n\to a$。
\end{proof}

\newpage
\begin{example}[数列 $\bigl(1+\frac1n\bigr)^n$]
令 $b_n=(1+1/n)^n$。则 $\{b_n\}$ 严格递增且有上界，记
\[
e:=\lim_{n\to\infty}\left(1+\frac1n\right)^n.
\]
\end{example}
\begin{proof}
二项式展开给出
\[
b_n=2+\sum_{k=2}^{n}\frac1{k!}
\prod_{j=1}^{k-1}\left(1-\frac jn\right).
\]
对固定的 $k\ge2$，当 $n$ 增加时，上式中的每个因子严格增加，
且 $b_{n+1}$ 比 $b_n$ 多一个正项，所以 $b_{n+1}>b_n$。
另一方面，对 $n\ge2$，
\[
b_n<2+\sum_{k=2}^{n}\frac1{k!}
\le2+\sum_{k=2}^{n}\frac1{k(k-1)}
=3-\frac1n<3.
\]
$b_1=2$ 也有上界 $3$，故由单调有界定理，极限 $e$ 存在。
\end{proof}

这个极限也可用于相应的实数列。若 $t_n>0$ 且 $t_n\to+\infty$，则
$\bigl(1+1/t_n\bigr)^{t_n}\to e$。事实上，取
$m_n=\lfloor t_n\rfloor$，当 $m_n\ge1$ 时有
\[
\left(1+\frac1{m_n+1}\right)^{m_n}
\le\left(1+\frac1{t_n}\right)^{t_n}
\le\left(1+\frac1{m_n}\right)^{m_n+1},
\]
两侧均趋于 $e$。由此，若 $u_n\to0$，且 $u_n>-1$、$u_n\ne0$，
则 $\bigl(1+u_n\bigr)^{1/u_n}\to e$。其中 $u_n>0$ 的项令
$t_n=1/u_n$；$u_n<0$ 的项令 $t_n=-1/u_n$，并用
\[
\left(1-\frac1{t_n}\right)^{-t_n}
=\left(1+\frac1{t_n-1}\right)^{t_n-1}
\left(1+\frac1{t_n-1}\right).
\]

\subsection*{子列与其他判别}
若 $n_1<n_2<\cdots$ 是严格递增的正整数列，则
$\{a_{n_k}\}_{k=1}^{\infty}$ 称为 $\{a_n\}$ 的一个\textbf{子列}。
由于 $n_k\ge k$，若 $a_n\to a$，则每个子列都有 $a_{n_k}\to a$。
因此，若能找到两个收敛于不同极限的子列，原数列就不收敛。
例如 $a_n=(-1)^n$ 的奇、偶子列分别收敛于 $-1$、$1$。

另外，对实数列还有如下判别：$\{a_n\}$ 收敛当且仅当它满足
\textbf{柯西条件}，即对任意 $\varepsilon>0$，存在 $N$，使
$m,n>N$ 时 $|a_m-a_n|<\varepsilon$。\textbf{魏尔斯特拉斯定理}
指出，每个有界实数列都含有收敛子列；这里要求原数列有界，
所得结论是某个子列收敛。
\end{document}
